\documentclass[11pt,a4paper]{article}

% ── packages ──────────────────────────────────────────────────────────────
\usepackage[utf8]{inputenc}
\usepackage[T1]{fontenc}
\usepackage{lmodern}
\usepackage[margin=1in]{geometry}
\usepackage{amsmath,amssymb,amsfonts}
\usepackage{graphicx}
\usepackage{booktabs}
\usepackage{multirow}
\usepackage{enumitem}
\usepackage{hyperref}
\usepackage{xcolor}
\usepackage{algorithm}
\usepackage{algpseudocode}
\usepackage{listings}
\usepackage{caption}
\usepackage{subcaption}
\usepackage{natbib}
\bibliographystyle{plainnat}

\hypersetup{colorlinks=true, linkcolor=blue!60!black, citecolor=blue!60!black, urlcolor=blue!60!black}

\lstset{
  basicstyle=\ttfamily\small,
  breaklines=true,
  frame=single,
  columns=fullflexible,
}

% ── title ─────────────────────────────────────────────────────────────────
\title{%
  \textbf{YtteriaDragon: A Sparse Mixture-of-Experts Language Model\\
  and Taglish Intent Corpus for On-Device\\
  Sari-Sari Store Intelligence}%
}

\author{
  GodYtteria\\
  Independent Researcher\\
  \texttt{github.com/GodYtteria}
}

\date{November 2026}

\begin{document}
\maketitle

% ══════════════════════════════════════════════════════════════════════════
\begin{abstract}
We present \textbf{YtteriaDragon}, a family of sparse Mixture-of-Experts (MoE)
language models designed for on-device deployment in resource-constrained
environments, and \textbf{Tinda360 Taglish Intent Corpus}, a 4{,}777-example
intent classification dataset for Filipino micro-retail (sari-sari store)
natural language understanding. The flagship \texttt{50M-tinda360} configuration
achieves 36.1M total parameters with only 4.8M active per token through
fine-grained MoE routing (32 experts, top-1), hash-based routing for early
layers following DeepSeek-V4, $\sqrt{\mathrm{softplus}}$ affinity scoring,
hybrid attention mixing full quadratic attention with Gated DeltaNet linear
attention, and a multi-token prediction (MTP) auxiliary head. The intent corpus
covers 23 store-operation intents in Taglish (Tagalog-English code-switching)
with systematic coverage of typographic noise, SMS abbreviations, emotional
registers, and disambiguation pairs. We describe the architecture design,
dataset construction methodology, and deployment considerations for
sub-100M models on mobile devices serving the Philippine sari-sari store
ecosystem of 1.3 million micro-retailers.
\end{abstract}

\textbf{Keywords:} sparse MoE, on-device NLP, Taglish, intent classification,
sari-sari store, Gated DeltaNet, hybrid attention, Filipino NLU

% ══════════════════════════════════════════════════════════════════════════
\section{Introduction}
\label{sec:intro}

The Philippines has approximately 1.3 million sari-sari stores---neighborhood
micro-retail shops that serve as the primary point of commerce for the majority
of Filipino households \citep{dti2023sari}. These stores are typically
owner-operated, with inventory management, credit tracking (``utang''), and
sales analysis performed manually or through basic mobile applications.

Recent advances in on-device language models have made it feasible to deploy
AI assistants directly on smartphones without cloud connectivity---a critical
requirement in areas with intermittent internet access. However, existing
models face two key challenges in this domain:

\begin{enumerate}[leftmargin=*]
  \item \textbf{Language gap.} Filipino store owners communicate in Taglish
    (Tagalog-English code-switching), SMS-style abbreviations, and regional
    dialects. No public intent classification corpus exists for this register
    in the micro-retail domain.
  \item \textbf{Compute budget.} On-device deployment on mid-range Android
    phones (4--6\,GB RAM, Snapdragon 600-series) requires models under 100M
    parameters with minimal active computation per token.
\end{enumerate}

We address both challenges with two contributions:

\begin{itemize}[leftmargin=*]
  \item \textbf{YtteriaDragon} (\S\ref{sec:architecture}): a sparse MoE
    architecture combining fine-grained expert routing (DeepSeek-V3 style),
    hash routing for early layers (DeepSeek-V4), $\sqrt{\mathrm{softplus}}$
    affinity scoring, hybrid quadratic/linear attention (Qwen3.5 style), and
    multi-token prediction. The \texttt{50M-tinda360} preset achieves 36.1M
    total / 4.8M active parameters.
  \item \textbf{Tinda360 Taglish Intent Corpus} (\S\ref{sec:dataset}): 4{,}777
    labeled intent examples across 23 store-operation intents, with a companion
    462-example training set, covering code-switching, typographic noise,
    emotional registers, and systematic disambiguation pairs.
\end{itemize}

% ══════════════════════════════════════════════════════════════════════════
\section{Related Work}
\label{sec:related}

\paragraph{Sparse Mixture-of-Experts.}
The MoE paradigm decouples total model capacity from per-token compute
\citep{shazeer2017outrageously}. Mixtral \citep{jiang2024mixtral} popularized
top-$k$ routing with softmax-renormalized weights. DeepSeek-V3
\citep{deepseekai2024deepseekv3} introduced \emph{fine-grained} MoE with
auxiliary-loss-free load balancing through per-expert bias terms, eliminating
the Switch Transformer's load-balancing auxiliary loss. DeepSeek-V4 extended
this with \emph{hash routing} for early layers where the routing signal from
hidden states is weak, using a fixed deterministic projection instead of a
learned gate.

\paragraph{Hybrid Attention.}
Qwen3.5 \citep{qwen2025qwen3} introduced a hybrid architecture alternating
between full quadratic attention (with gated queries, per-head RMS normalization,
and partial RoPE) and Gated DeltaNet linear attention
\citep{yang2024gated}. This achieves $O(T)$ recurrent inference with periodic
full-attention ``anchor'' layers for global context. The gated delta rule
extends classical delta-rule linear attention with a learned decay gate,
enabling selective memory updates.

\paragraph{Multi-Token Prediction.}
Meta \citep{gloeckle2024better} demonstrated that training language models to
predict multiple future tokens simultaneously (via parallel prediction heads)
improves sample efficiency and downstream task performance. The auxiliary MTP
heads are discarded at inference, preserving the base model's architecture.

\paragraph{Filipino NLP.}
Filipino NLP resources remain scarce. ParaCrawl and CC-100 include some
Tagalog web text, and \citet{cruz2020establishing} established benchmarks
for Filipino text classification. However, no public dataset addresses
Taglish intent classification in the micro-retail domain, where code-switching
density and SMS abbreviation rates are significantly higher than in formal text.

% ══════════════════════════════════════════════════════════════════════════
\section{Architecture}
\label{sec:architecture}

YtteriaDragon is a decoder-only transformer with the following configurable
components. Table~\ref{tab:presets} lists the full scaling ladder.

\subsection{Base Transformer}

The backbone follows the LLaMA architecture \citep{touvron2023llama}: RMSNorm
pre-normalization, SwiGLU feed-forward networks, Grouped-Query Attention (GQA)
with RoPE positional embeddings, and tied input/output embeddings. We use the
``rotate\_half'' RoPE layout for compatibility with the GGUF quantization
format.

\subsection{Fine-Grained Sparse MoE}
\label{sec:fgmoe}

Each transformer layer's FFN is replaced by a fine-grained MoE block containing
$E$ routed SwiGLU experts and one \emph{shared expert} that processes every
token (following DeepSeek-V3). The routing mechanism selects top-$k$ experts
per token:

\begin{enumerate}[leftmargin=*]
  \item A learned linear gate computes logits $\ell_i = W_g \cdot h$ for each
    expert $i \in \{1, \ldots, E\}$.
  \item Per-expert bias terms $b_i$ (not learned by gradient descent) are added
    to the routing logits: $r_i = \ell_i + b_i$.
  \item The top-$k$ experts by $r_i$ are selected.
  \item Output weights are computed from the \emph{original} logits (not biased)
    using the affinity function (\S\ref{sec:affinity}).
  \item Bias terms are updated each training step:
    $b_i \leftarrow b_i - \alpha \cdot (f_i - 1/E)$, where $f_i$ is the
    routed fraction and $\alpha$ is the bias learning rate.
\end{enumerate}

This eliminates the auxiliary load-balancing loss entirely while maintaining
balanced expert utilization.

\subsubsection{Hash Routing for Early Layers}
\label{sec:hash}

Following DeepSeek-V4, the first $H$ layers use \emph{hash routing}: a fixed
(non-learnable) random projection $P \in \mathbb{R}^{d \times E}$ replaces the
learned gate. Expert selection is determined by $\arg\text{topk}(h \cdot P, k)$
with uniform output weights ($w_i = 1/k$). This provides:

\begin{itemize}[leftmargin=*]
  \item \textbf{Perfect load balance} without bias correction.
  \item \textbf{Zero routing compute} (the projection is a buffer, not a parameter).
  \item \textbf{Stability} in early layers where hidden-state routing signals
    are noisy and unreliable.
\end{itemize}

The random projection is initialized once at model creation and registered as a
non-persistent buffer.

\subsubsection{$\sqrt{\mathrm{softplus}}$ Affinity Scoring}
\label{sec:affinity}

For non-hash layers, we replace the standard softmax affinity with
$\sqrt{\mathrm{softplus}}$ scoring:

\begin{equation}
  w_i = \frac{\sqrt{\mathrm{softplus}(\ell_i)}}
             {\sum_{j \in \text{top-}k} \sqrt{\mathrm{softplus}(\ell_j)}}
\end{equation}

where $\mathrm{softplus}(x) = \ln(1 + e^x)$. This provides smoother gradients
near zero compared to softmax and avoids the ``winner-take-all'' saturation
that softmax exhibits with large logit gaps. For top-1 routing (as in the
\texttt{50M-tinda360} preset), the weight is always 1.0 regardless of scoring
function; the benefit emerges in presets with top-$k > 1$.

\subsection{Hybrid Attention}
\label{sec:hybrid}

When \texttt{hybrid}$= N > 0$, every $N$-th layer uses full quadratic attention
(GatedAttention) and the remaining layers use Gated DeltaNet linear attention:

\paragraph{GatedAttention (full layers).}
Following Qwen3.5: the query projection is 2$\times$ wide, producing both
queries and a sigmoid output gate. Per-head RMSNorm is applied to Q and K.
Partial RoPE is applied to the first \texttt{rope\_frac} of the head dimension.

\paragraph{GatedDeltaNet (linear layers).}
The gated delta-rule recurrence \citep{yang2024gated} with:
(a) a 1-D causal convolution on the Q/K/V projection,
(b) learned decay gate $g$ and update gate $\beta$,
(c) L2-normalized Q/K,
(d) chunked $O(T)$ recurrence with 64-token chunks.
At inference, the recurrence state $S \in \mathbb{R}^{H \times d_k \times d_v}$
replaces the KV cache, giving constant memory regardless of sequence length.

\paragraph{Attention pattern for \texttt{50M-tinda360}.}
With 8 layers and \texttt{hybrid}$=4$: layers 3 and 7 are GatedAttention;
layers 0, 1, 2, 4, 5, 6 are GatedDeltaNet. This gives 75\% linear attention
(efficient inference) with 25\% full attention (global context anchors).

\subsection{Multi-Token Prediction}
\label{sec:mtp}

Following \citet{gloeckle2024better}, $M$ auxiliary prediction heads are added
during training. Head $k$ consists of one additional transformer block applied
to the trunk output, followed by the shared final norm and language model head.
Head $k$ is trained to predict token $t + 2 + k$ given position $t$. The MTP
loss is:

\begin{equation}
  \mathcal{L}_{\text{MTP}} = \frac{1}{M} \sum_{k=0}^{M-1}
    \mathrm{CE}(\text{head}_k(h_t), y_{t+2+k})
\end{equation}

weighted by $\lambda_{\text{MTP}} = 0.3$. The heads are discarded at inference
(GGUF export drops them), so the deployed model is unchanged.

\subsection{Model Configurations}

\begin{table}[h]
\centering
\caption{YtteriaDragon scaling ladder. Active params count top-$k$ routed
experts + shared expert per MoE layer. MTP heads excluded from active count
(inference-only metric). The \texttt{50M-tinda360} preset is the focus of
this paper.}
\label{tab:presets}
\small
\begin{tabular}{lrrrrrrr}
\toprule
\textbf{Preset} & \textbf{$d$} & \textbf{Layers} & \textbf{Experts} &
  \textbf{Top-$k$} & \textbf{Hybrid} & \textbf{Total} & \textbf{Active} \\
\midrule
10M         & 256  & 4  & --  & -- & -- & 10.0M & 10.0M \\
30M         & 384  & 10 & --  & -- & -- & 29.5M & 29.5M \\
30M-moe8    & 384  & 10 & 8   & 2  & -- & 63.7M & 28.4M \\
30M-fgmoe   & 384  & 10 & 32  & 2  & -- & \multicolumn{2}{c}{(fg, shared)} \\
30M-hybrid  & 384  & 10 & --  & -- & 4  & 28.5M & 28.5M \\
\textbf{50M-tinda360} & \textbf{192} & \textbf{8} & \textbf{32} &
  \textbf{1} & \textbf{4} & \textbf{36.1M} & \textbf{4.8M} \\
100M        & 640  & 18 & --  & -- & -- & 99.8M & 99.8M \\
300M        & 1024 & 24 & --  & -- & -- & 303M  & 303M \\
500M        & 1280 & 26 & --  & -- & -- & 509M  & 509M \\
\bottomrule
\end{tabular}
\end{table}

\paragraph{\texttt{50M-tinda360} configuration.}
\texttt{vocab\_size}$=8192$, \texttt{max\_seq\_len}$=1024$, $d_{\text{model}}=192$,
8 layers, 4 attention heads, 2 KV heads, 32 FG experts with
\texttt{expert\_ffn}$=192$, top-1 routing, hash routing for layers 0--1,
$\sqrt{\mathrm{softplus}}$ affinity, \texttt{hybrid}$=4$, MTP$=1$,
\texttt{attn\_window}$=512$.

The design rationale is \emph{capacity without compute}: 32 experts provide
7.5$\times$ the FFN capacity of a dense model at the same active parameter
count, enabling the model to learn specialized expert pathways for each of
the 23 intents while keeping inference fast enough for real-time mobile use.

\subsection{Auto-Compaction and Sliding Window}
\label{sec:compact}

The hybrid architecture provides built-in context compaction:

\begin{itemize}[leftmargin=*]
  \item \textbf{GatedDeltaNet layers} (6 of 8): $O(1)$ memory regardless of
    context length. The recurrence state $S \in \mathbb{R}^{H \times d_k \times d_v}$
    is a fixed-size lossy compression of the full history ($\sim$37\,KB per layer).
  \item \textbf{GatedAttention layers} (2 of 8): sliding window KV eviction with
    \texttt{attn\_window}$=512$ tokens. Older entries are evicted automatically,
    bounding memory to $\sim$384\,KB per layer.
\end{itemize}

Total inference memory at 1024-token context: $\sim$1\,MB (220\,KB GDN states +
768\,KB KV cache). An explicit \texttt{auto\_compact()} method can be called
between user turns to proactively evict stale KV entries during idle time.

\subsection{Tool Calling}
\label{sec:tools}

The model supports structured tool calling via a ChatML template with special
tokens (\texttt{<tool\_call>}, \texttt{</tool\_call>}, \texttt{<tool\_result>},
\texttt{</tool\_result>}). Twenty store-operation functions are defined
(e.g., \texttt{get\_sales\_summary}, \texttt{get\_product\_info},
\texttt{get\_utang\_list}). The training data includes:

\begin{itemize}[leftmargin=*]
  \item \textbf{Single-call}: user query $\to$ one tool call $\to$ natural reply.
  \item \textbf{Parallel multi-call}: query requiring 2--3 simultaneous tool calls.
  \item \textbf{Multi-turn}: conversational sequences with follow-up questions.
  \item \textbf{Abstention}: queries outside the tool scope, trained to decline.
  \item \textbf{Function masking}: 33\% of examples use randomized tool names
    to prevent memorization.
\end{itemize}

The model is trained to never generate data it did not receive from a tool call,
preventing hallucination of store metrics.

\subsection{Training Infrastructure}

\paragraph{Optimizer.}
We use the Muon optimizer \citep{jordan2024muon} for 2-D weight matrices
(attention projections, expert weights) with Newton-Schulz orthogonalization,
paired with AdamW for embeddings, norms, and routing parameters. This follows
the Qwen3.8-Flash-Next training recipe. Learning rate follows a cosine schedule
decaying to $1/10$ of peak (consensus from DeepSeek-V4, Arcee Trinity, IMU-1),
with 1000-step linear warmup.

\paragraph{Stability techniques.}
SwiGLU activation clamping (DeepSeek-V4), router z-loss (ST-MoE),
MoE-specific truncated normal initialization, and BF16 mixed-precision
training with the router forced to FP32 (nanoMoE best practice).

\paragraph{Training data.}
The combined training corpus consists of 12{,}426 examples: 9{,}938 intent
classification, 1{,}594 simplified training, and 894 tool-calling sequences.

% ══════════════════════════════════════════════════════════════════════════
\section{Tinda360 Taglish Intent Corpus}
\label{sec:dataset}

\subsection{Overview}

The corpus consists of three files:

\begin{itemize}[leftmargin=*]
  \item \textbf{Intent Classifier} (9{,}938 examples): structured JSON output
    with intent label and optional entity/timeframe fields.
  \item \textbf{Training Data} (1{,}594 examples): simplified format with bare
    intent labels, suitable for lightweight fine-tuning.
  \item \textbf{Tool Calling} (894 examples): multi-turn ChatML conversations
    with tool call/result sequences, including 207 function-masked and 33
    abstention examples.
\end{itemize}

All three files cover 23 intents and 20 tool functions spanning the full
operational scope of a sari-sari store AI assistant.

\subsection{Intent Taxonomy}

\begin{table}[h]
\centering
\caption{Intent taxonomy with example counts (classifier / training data).}
\label{tab:intents}
\small
\begin{tabular}{llrr}
\toprule
\textbf{Intent} & \textbf{Description} & \textbf{Clf.} & \textbf{TD} \\
\midrule
\texttt{product\_lookup}    & Price / stock check for a product   & 2{,}332 & 22 \\
\texttt{stock\_forecast}    & Predict when product runs out       & 381  & 21 \\
\texttt{sales}              & Sales/revenue report by timeframe   & 212  & 28 \\
\texttt{greeting}           & Greeting or conversation opener     & 133  & 29 \\
\texttt{top\_products}      & Best-selling products               & 112  & 20 \\
\texttt{low\_stock}         & Items at critically low stock       & 110  & 22 \\
\texttt{help}               & How to use the assistant            & 110  & 24 \\
\texttt{utang}              & Credit/debt tracking                & 109  & 23 \\
\texttt{overview}           & General store status summary        & 102  & 19 \\
\texttt{sales\_compare}     & Compare sales across periods        & 101  & 20 \\
\texttt{expiryAlert}        & Expiring / expired products         & 101  & 21 \\
\texttt{thanks}             & Gratitude / conversation closer     & 100  & 23 \\
\texttt{expenses}           & Operating cost report               & 100  & 17 \\
\texttt{slow\_moving}       & Non-selling / stagnant inventory    & 86   & 18 \\
\texttt{restock\_budget}    & Reorder suggestions within budget   & 85   & 19 \\
\texttt{utang\_risk}        & High-risk debtors / bad debts       & 83   & 19 \\
\texttt{suki\_loyalty}      & Loyal / frequent customer report    & 83   & 18 \\
\texttt{payment\_breakdown} & Cash vs.\ GCash vs.\ credit split  & 80   & 18 \\
\texttt{profit\_margin}     & Markup and profitability analysis   & 82   & 17 \\
\texttt{sales\_trend}       & Sales trajectory over time          & 79   & 19 \\
\texttt{busiest\_hours}     & Peak selling hours / days           & 75   & 17 \\
\texttt{product\_count}     & Total inventory item count          & 63   & 14 \\
\texttt{store\_info}        & Store profile / registration info   & 58   & 14 \\
\midrule
\textbf{Total}              &                                     & \textbf{4{,}777} & \textbf{462} \\
\bottomrule
\end{tabular}
\end{table}

\subsection{Linguistic Diversity}
\label{sec:diversity}

The corpus systematically covers the following variation axes:

\paragraph{Code-switching density.}
Examples range from pure Tagalog (``Magkano po ang bigas'') to pure English
(``How much is rice'') with dense Taglish mixing (``Check mo yung stock ng
softdrinks'').

\paragraph{SMS / text-speak abbreviations.}
Common Filipino texting patterns: vowel dropping (``Mgkno'' $\to$ ``Magkano''),
consonant clusters (``Stck'' $\to$ ``Stock''), phonetic spelling
(``Tenkyu'' $\to$ ``Thank you'').

\paragraph{Typographic noise.}
ALL CAPS (``MAGKANO BIGAS''), excessive punctuation (``Sales?????''),
no punctuation, mixed case.

\paragraph{Register variation.}
Polite (``po'', ``naman'', ``please'', ``paki-''), casual (``Uy'', ``Oi''),
Filipino internet slang (``Lodi'', ``Bes'', ``Charot''), emotional / frustrated
(``Naku ang daming utang!'', ``Bakit ang baba ng benta!'').

\paragraph{Disambiguation pairs.}
40+ systematically constructed pairs that test routing accuracy on ambiguous
stems. For example, ``Kumusta'' $\to$ \texttt{greeting} vs.\
``Kumusta stock'' $\to$ \texttt{low\_stock} vs.\
``Kumusta benta'' $\to$ \texttt{sales}.

\paragraph{Edge cases.}
Gibberish/noise (``asdf'', ``hahaha'' $\to$ \texttt{greeting} fallback),
single-word queries (``Utang'' $\to$ \texttt{utang}),
compound queries (``Sales at expenses ngayon'' $\to$ primary intent),
negation (``Wala bang benta?'' $\to$ \texttt{sales}).

\subsection{Entity and Timeframe Annotation}

For \texttt{product\_lookup} and \texttt{stock\_forecast} intents, the output
includes an \texttt{entity} field with the normalized product name. For
time-sensitive intents (\texttt{sales}, \texttt{top\_products},
\texttt{expenses}), a \texttt{timeframe} field (``today'', ``week'',
``month'') is included. This enables the downstream application to issue
structured database queries directly from the classifier output.

\subsection{Product Coverage}

The corpus includes 75+ generic sari-sari store product categories (``bigas'',
``itlog'', ``sardinas'', ``mantika'', etc.) representing the typical inventory
of a Filipino neighborhood store. Product names are kept generic (not
brand-specific) to improve generalization across stores with different
supplier relationships.

% ══════════════════════════════════════════════════════════════════════════
\section{Deployment Considerations}
\label{sec:deployment}

\subsection{Target Platform}

Tinda360 is a Flutter-based mobile application targeting Android devices in
the Philippine market. The AI assistant (``Yumi'') runs entirely on-device
to support offline operation in areas with limited connectivity.

\subsection{Quantization and Inference}

The YtteriaDragon architecture is designed for GGUF export and llama.cpp
inference:

\begin{itemize}[leftmargin=*]
  \item \textbf{Tied embeddings} for smaller model files.
  \item \textbf{FFN dimensions as multiples of 256} to ensure k-quant
    quantization (Q4\_K\_M) works without fallback to lower-quality formats.
  \item \textbf{MTP heads dropped} at export time (inference model unchanged).
  \item \textbf{Hybrid attention exports as Qwen3\_5ForCausalLM} for
    llama.cpp compatibility (QWEN35 architecture tag).
\end{itemize}

At Q4\_K\_M quantization, the \texttt{50M-tinda360} model is estimated at
$\sim$18\,MB, well within the memory constraints of mid-range Android devices.

\subsection{Inference Characteristics}

\begin{itemize}[leftmargin=*]
  \item \textbf{Linear attention layers}: $O(1)$ per token (recurrence state
    instead of KV cache), enabling fast autoregressive generation.
  \item \textbf{Full attention anchors}: every 4th layer maintains a standard
    KV cache, but \texttt{max\_seq\_len}$=512$ keeps this small.
  \item \textbf{Top-1 routing}: only 1 of 32 experts executes per layer,
    minimizing memory bandwidth for expert weight loading.
  \item \textbf{Active parameters}: 4.8M $\approx$ 2.4\,MB at FP16, fitting
    entirely in L2 cache on modern mobile SoCs.
\end{itemize}

% ══════════════════════════════════════════════════════════════════════════
\section{Limitations and Future Work}
\label{sec:limitations}

\begin{enumerate}[leftmargin=*]
  \item \textbf{No training results.} This paper describes the architecture
    and dataset; training, evaluation metrics (accuracy, latency, perplexity),
    and ablation studies will be reported in a future revision once GPU
    compute is available.
  \item \textbf{Single annotator.} The intent corpus was constructed by a
    single developer with native Filipino fluency. Inter-annotator agreement
    scores and crowd-sourced validation are planned.
  \item \textbf{Limited regional coverage.} The corpus is Tagalog-English
    dominant. Bisaya, Ilokano, and other major Philippine languages are
    underrepresented.
  \item \textbf{No real user data.} All examples are synthetic. Field testing
    with actual sari-sari store owners will be necessary to validate coverage
    and discover missing intents.
  \item \textbf{Tokenizer.} A BPE tokenizer trained on the Taglish corpus
    has not yet been produced. The vocabulary size (8{,}192) is chosen to
    balance coverage with embedding table size.
\end{enumerate}

% ══════════════════════════════════════════════════════════════════════════
\section{Conclusion}
\label{sec:conclusion}

We presented YtteriaDragon, a sparse MoE language model architecture
incorporating hash routing, $\sqrt{\mathrm{softplus}}$ affinity scoring,
hybrid quadratic/linear attention, and multi-token prediction, targeting
on-device deployment for Filipino sari-sari store intelligence. The
\texttt{50M-tinda360} configuration achieves a 7.5$\times$ capacity-to-compute
ratio (36.1M total / 4.8M active parameters) through fine-grained top-1
routing over 32 experts.

The accompanying Tinda360 Taglish Intent Corpus provides 4{,}777 labeled
examples across 23 store-operation intents, systematically covering the
code-switching, typographic, and register variation characteristic of Filipino
micro-retail communication. Together, these contributions aim to bring
AI-powered business intelligence to the 1.3 million sari-sari stores that
form the backbone of Philippine neighborhood commerce.

The model architecture, dataset, and training code are open-source at
\url{https://github.com/GodYtteria/Tinda360}.

% ══════════════════════════════════════════════════════════════════════════
\begin{thebibliography}{20}

\bibitem[Cruz and Cheng(2020)]{cruz2020establishing}
Cruz, J.~C.~B. and Cheng, C.
\newblock Establishing baselines for text classification in low-resource
  languages.
\newblock \emph{arXiv preprint arXiv:2005.02068}, 2020.

\bibitem[{DeepSeek-AI}(2024)]{deepseekai2024deepseekv3}
{DeepSeek-AI}.
\newblock DeepSeek-V3 technical report.
\newblock \emph{arXiv preprint arXiv:2412.19437}, 2024.

\bibitem[{DTI}(2023)]{dti2023sari}
{Department of Trade and Industry, Philippines}.
\newblock Sari-sari store development program annual report, 2023.

\bibitem[Gloeckle et~al.(2024)]{gloeckle2024better}
Gloeckle, F., Idrissi, B.~Y., Rozière, B., Lopez-Paz, D., and Synnaeve, G.
\newblock Better \& faster large language models via multi-token prediction.
\newblock \emph{arXiv preprint arXiv:2404.19737}, 2024.

\bibitem[Jiang et~al.(2024)]{jiang2024mixtral}
Jiang, A.~Q., Sablayrolles, A., Roux, A., Mensch, A., et~al.
\newblock Mixtral of experts.
\newblock \emph{arXiv preprint arXiv:2401.04088}, 2024.

\bibitem[Jordan(2024)]{jordan2024muon}
Jordan, K.
\newblock Muon: An optimizer for hidden layers in neural networks.
\newblock \emph{Blog post}, 2024.

\bibitem[{Qwen Team}(2025)]{qwen2025qwen3}
{Qwen Team}.
\newblock Qwen3 technical report.
\newblock \emph{arXiv preprint arXiv:2505.09388}, 2025.

\bibitem[Shazeer et~al.(2017)]{shazeer2017outrageously}
Shazeer, N., Mirhoseini, A., Maziarz, K., et~al.
\newblock Outrageously large neural networks: The sparsely-gated
  mixture-of-experts layer.
\newblock \emph{ICLR}, 2017.

\bibitem[Touvron et~al.(2023)]{touvron2023llama}
Touvron, H., Lavril, T., Izacard, G., et~al.
\newblock LLaMA: Open and efficient foundation language models.
\newblock \emph{arXiv preprint arXiv:2302.13971}, 2023.

\bibitem[Yang et~al.(2024)]{yang2024gated}
Yang, S., Wang, B., Shen, Y., Panda, R., and Kim, Y.
\newblock Gated delta networks: Improving mamba2 with delta rule.
\newblock \emph{arXiv preprint arXiv:2412.06464}, 2024.

\end{thebibliography}

\end{document}
